\subsection{射影线性映射}

设 V、 \(V'\) 是域 K 上的两个左向量空间，f 是 V 到 \(V'\) 的一个线性映射， \(N = \bar{f}^{-1}(0)\) 是它的核。立即可见，V 中一条不含于 N 的（过 0 的）直线在 f 下的像是 \(V'\) 中一条（过 0 的）直线；因此，在过渡到商时，f 定义了一个映射 g （从 \(\mathbf{P}(\mathbf{V}) - \mathbf{P}(\mathbf{N})\) 到 \(\mathbf{P}(\mathbf{V}')\) ）。这样的映射称为射影线性映射（或简称射影映射）；尽管它定义在 \(\mathbf{P}(\mathbf{V}) - \mathbf{P}(\mathbf{N})\) 上而不是 \(\mathbf{P}(\mathbf{V})\) 上（当 \(N \neq \{0\}\) 时），我们仍按语言滥用的说法称 g 是 \(\mathbf{P}(\mathbf{V})\) 到 \(\mathbf{P}(\mathbf{V}')\) 的一个射影映射。射影线性簇 \(\mathbf{P}(\mathbf{N})\) （g 在其上无定义）称为 g 的中心。

注意，当 \(g\) 定义在整个 \(\mathbf{P}(\mathbf{V})\) 上时（即当 \(\mathbf{N} = \{0\}\) 时）， \(g\) 是 \(\mathbf{P}(\mathbf{V})\) 到 \(\mathbf{P}(\mathbf{V}')\) 的一个单射。

当基 \((a_{\lambda})_{\lambda\in\mathbf{L}}, (b_{\mu})_{\mu\in\mathbf{M}}\) 分别在 V 与 \(V'\) 中给定时， \(\mathbf{P}(\mathbf{V})\) 到 \(\mathbf{P}(\mathbf{V}')\) 的一个射影映射把 \(\mathbf{P}(\mathbf{V})\) 中具有齐次坐标 \(\xi_{\lambda}\) ( \(\lambda \in \mathbf{L}\) ) 的点映到 \(\mathbf{P}(\mathbf{V}')\) 中的一点，该点具有形如以下的齐次坐标系 \(\eta_{\mu}\) ( \(\mu \in \mathbf{M}\) )

\[
\eta_ {\mu} = \sum_ {\lambda \in L} \xi_ {\lambda} \alpha_ {\lambda \mu} \quad (\alpha_ {\lambda \mu} \in K).\tag{6}
\]

 \(g\) 的中心是由以下方程定义的线性簇

\[
\sum_ {\lambda \in L} \xi_ {\lambda} \alpha_ {\lambda \mu} = 0 \quad (\mu \in M).
\]

如果 C 是 g 的中心，且 M 是 \(\mathbf{P}(\mathbf{V})\) 的一个线性簇，则 \(\mathbf{M} - (\mathbf{M} \cap \mathbf{C})\) 在 g 下的像是 \(\mathbf{P}(\mathbf{V}')\) 的一个线性簇，（按语言滥用）记作 \(g(\mathbf{M})\) 。则

\[
\dim g (\mathbf {M}) + \dim (\mathbf {M} \cap \mathbf {C}) + 1 = \dim \mathbf {M}\tag{7}
\]

（§7，第 4 号，公式 (12)）。若 \(\mathbf{M}'\) 是 \(\mathbf{P}(\mathbf{V}')\) 的一个线性簇，则 \(\bar{g}^{-1}(\mathbf{M}') \cup \mathbf{C}\) 是 \(\mathbf{P}(\mathbf{V})\) 的一个线性簇，并且

\[
\dim (\bar {g} ^ {- 1} (\mathbf {M} ^ {\prime}) \cup \mathbf {C}) = \dim \mathbf {C} + \dim (\mathbf {M} ^ {\prime} \cap g (\mathbf {P} (\mathbf {V}))) + 1.\tag{8}
\]

我们按语言滥用的说法，称 \(\bar{g}^{-1}(\mathbf{M}') \cup \mathbf{C}\) 为 \(\mathbf{M}'\) 在 \(g\) 下的逆像。

由于线性映射在基 \((e_{t})\) （属于 V ）上取的值可以在 \(V'\) 中任意选取，可见存在 \(\mathbf{P}(\mathbf{V})\) 到 \(\mathbf{P}(\mathbf{V}')\) 的一个射影线性映射，它在点 \(\pi(e_{t})\) 处取任意的值。但是（即使当 g 处处有定义时）给出 \(g(\pi(e_{t}))\) 也不能唯一地确定 g（习题 10）。

两个双射的射影映射的复合仍是一个射影映射；这样的双射的逆映射也是射影映射。射影空间 \(\mathbf{P}(\mathbf{V})\) 到自身上的双射射影映射因此构成一个群，称为 \(\mathbf{P}(\mathbf{V})\) 的射影群，记作 \(\mathbf{PGL}(\mathbf{V})\) ；我们用 \(\mathbf{PGL}_{n}(\mathbf{K})\) 或 \(\mathbf{PGL}(n,\mathbf{K})\) 代替 \(\mathbf{PGL}(\mathbf{K}_{s}^{n})\) 。

注记。在射影空间 \(\mathbf{P}(\mathbf{V})\) （在域 \(\mathbf{K}\) 上）中，设 \(\mathbf{H} = \mathbf{P}(\mathbf{W})\) 是一个超平面。存在一个双射线性映射 \(f\) （从 \(\mathbf{V}\) 映到 \(\mathbf{K}_s \times \mathbf{W}\) 上），使得 \(f(\mathbf{W}) = \mathbf{W}\) ；设 \(g\) 是由 \(f\) 通过过渡到商而得到的射影映射。已经看到（第 8 号）， \(\mathbf{P}(\mathbf{W})\) 在 \(\mathbf{P}(\mathbf{K}_s \times \mathbf{W})\) 中的补集可以等同于一个仿射空间，其平移空间是 \(\mathbf{W}\) 。当 \(\mathbf{P}(\mathbf{V})\) 与 \(\mathbf{P}(\mathbf{K}_s \times \mathbf{W})\) 借助 \(g\) 等同时，就说取 \(\mathbf{H}\) 为 \(\mathbf{P}(\mathbf{V})\) 中的无穷远超平面；于是 \(\mathbf{H}\) 在 \(\mathbf{P}(\mathbf{V})\) 中的补集等同于一个仿射空间，其平移空间是 \(\mathbf{W}\) 。

